% Section 11.8, from lectures/L11/appendix/b_exercises.md.

\section{Exercises}
\label{c11:sec:exercises}

\begin{aside}
These exercises consolidate \secref{c11:sec:arch} to \ref{c11:sec:regmap}: the block architecture,
the interface it lands in, and the register map it maps onto. The \code{can_def} package they all
rest on was written in \chapref{c10:ch}.

Exercise~\ref{c11:ex:reason} and \ref{c11:ex:groups} are done on paper. Exercise~\ref{c11:ex:entity}
and \ref{c11:ex:iface} are done at the keyboard, against the
\file{controller/can_controller.vhd} written during this session. \code{meta_prev} and its exercises
belong to \chapref{c12:ch}, together with \code{bit_timer}.
\end{aside}

% ------------------------------------------------------------------------------------------
\exerciseset{Reasoning about the architecture}

\exercise{Who owns what}{Reasoning}
\label{c11:ex:reason}
\xpart{a} \code{tx_shift_reg} and \code{rx_shift_reg} do not know whether the bits they shift belong
to the ID, the control field, a data byte or the CRC value. Explain, with reference to the block
diagram, \emph{which} block does know that, and how it can reuse the same shift register for all of
it instead of needing five different ones.

\xpart{b} \code{crc15} is fed one bit at a time and told when to stop by whichever of
\code{tx_shift_reg} and \code{rx_shift_reg} is active at the time (via their outputs \code{stuff} and
\code{real_bit_valid} respectively, which gate \code{crc15.enable}). Why does \code{crc15} have to be
told to \emph{skip} stuff bits in particular, instead of simply being fed every bit that turns up on
the wire?

\xpart{c} \code{can_controller} instantiates the four other blocks inside itself, so from outside the
whole controller is a single component. Suppose someone proposes factoring the arbitration-loss check
out into a block of its own \emph{beside} \code{can_controller}. Explain what that block would need
access to, and why that makes the solution worse than leaving the check where it is.

\exercise{Splitting a different frame into groups}{Reasoning}
\label{c11:ex:groups}
\Secref{c11:sec:trace} follows a frame with DLC \code{2} through the blocks and arrives at it
reaching the bus as eight loads of \code{tx_shift_reg}, 50 bits together, plus a 10-bit tail driven
directly.

\xpart{a} Redo that split for a frame with ID \code{0x7A0} and DLC \code{5}. State the list of groups
(every group's contents and bit count, in transmission order), the total number of bits in the
stuffed region, and the total number of bits in the frame before any stuff bit is inserted.

\xpart{b} How many groups would a frame with DLC \code{0} need, and which entries from your list in
\textbf{a)} disappear? What does that tell you about how \code{can_controller} has to be structured
around the data field, compared with the identifier or the CRC?

\xpart{c} Four of the groups in your list are narrower than 8 bits, and \code{tx_shift_reg}'s data
input is always a whole byte. Name them, and say what \code{can_controller} has to pass along with
the data for the register to send the right number of bits. (You do not need \chapref{c14:ch}'s
answer to \emph{how}; the question is what information the register cannot work out on its own.)

\xpart{d} None of your group counts is the number of bits actually driven out on the wire. Explain
what is missing, why it cannot be known from the DLC alone, and which block is responsible for it.

% ------------------------------------------------------------------------------------------
\exerciseset{The entity and the interface}

\exercise{The entity, checked}{Simulation}
\label{c11:ex:entity}
\xpart{a} Run \code{make build-project} from the root of the group's repository, exactly as in
\chapref{c10:ch}. Two of your own files should now be reported, alongside the two \file{bridge/}
files handed out, and every testbench should still be skipped:

\begin{console}
--> controller/can_def.vhd analyzes cleanly
--> controller/can_controller.vhd analyzes cleanly
--> bridge/spi_def.vhd analyzes cleanly
--> bridge/spi_slave.vhd analyzes cleanly

==> meta_prev_tb (skipped: no meta_prev.vhd yet)
==> bit_timer_tb (skipped: no bit_timer.vhd yet)
==> crc15_tb (skipped: no crc15.vhd yet)
==> tx_shift_reg_tb (skipped: no tx_shift_reg.vhd yet)
==> rx_shift_reg_tb (skipped: no rx_shift_reg.vhd yet)
==> can_controller_tb (skipped: no meta_prev.vhd bit_timer.vhd crc15.vhd tx_shift_reg.vhd rx_shift_reg.vhd yet)
==> register_bank_tb (skipped: no register_bank.vhd yet)
==> spi_reg_bridge_tb (skipped: no spi_reg_bridge.vhd yet)

0 testbench(es) run, 8 skipped.
\end{console}

The second line is the whole point of this exercise: an entity with an empty architecture is still a
complete, valid design unit. Now go beyond what \code{make build-project} does and elaborate it by
hand, which is the second of GHDL's three steps:

\begin{shell}
cd controller
ghdl -a --std=93 can_def.vhd can_controller.vhd
ghdl -e --std=93 can_controller
\end{shell}

It elaborates cleanly, with no subblocks inside it and nothing to simulate. There is no third step to
run here, since a design with no behaviour has nothing to check and no testbench of its own;
\code{can_controller_tb} drives the \emph{finished} controller and is waiting for \chapref{c17:ch}.
It is in \chapref{c12:ch} that the third step turns up, with the first module that has a testbench of
its own.

\xpart{b} \code{can_controller_tb} is already in \file{controller/}, and it needs seven modules, two
of which you have written. Say which five are missing and in which chapter each one arrives. Why is
it better for the build to skip that testbench than for it to try and fail?

\xpart{c} \emph{Optional, and only if you have Quartus Prime Lite installed.} Take the same file
through Quartus, following appendix~\ref{app:quartus}: New Project Wizard, device
\code{5CEBA4F23C7N}, add \file{can_def.vhd} and \file{can_controller.vhd}, set
\code{can_controller} as the top-level entity, and compile.

It passes. Read the warnings rather than skipping them, and say which of them are unavoidable for an
architecture that drives nothing at all. Then find the pin count in the fitter's report and compare
it with what a DE0-CV actually offers a human being: ten switches, four buttons, ten LEDs. Write down
what you think has to sit between this entity and the board, and keep the note; \chapref{c18:ch} and
\ref{c19:ch} build exactly that, and call it \code{register_bank} and \code{can_spi_node}.

Nothing later in the book depends on this, and \chapref{c12:ch} to \ref{c19:ch} need only GHDL. It is
here because that pin count alone is the whole argument for the two layers above the controller,
seven chapters in advance, and reading it off a real fitter report makes it yours rather than a
number the book asserts.

\xpart{d} Temporarily swap \code{tx_bus} and \code{bus_en} in the entity's port list, so that they
stand in reverse order, and rerun \code{make build-project}. Both are outputs of type
\code{std_logic}, so the file still analyses cleanly. Explain why nothing catches it now, in which
chapter it would first do damage, and what that damage would look like on the bus (the clue is
\secref{c11:sec:interface}'s section on the bus side's five ports). Then say why a mix-up between two
ports of the same type is harder to find than a misspelled name would have been, given that
everything in this book is bound positionally. Swap them back.

\exercise{Reading the interface}{Reasoning}
\label{c11:ex:iface}
\xpart{a} A caller wants to send a frame with ID \code{0x123}, DLC \code{3} and the data bytes
\code{AA BB CC}. List every port it has to drive, and the value of each. \code{tx_data} is
\code{data_t}, 64 bits wide, and byte 0 sits in the most significant bits: give it in full.

\xpart{b} The same caller now wants to know whether the frame was sent. It has three status ports at
its disposal (\code{tx_done}, \code{rx_valid}, \code{error}). Which does it poll, and what is the one
thing it can \emph{never} find out from them, however often it polls? (The clue is
\secref{c11:sec:interface}'s remark about what \code{error} means; the reason is
\secref{c10:sec:frameformat}'s ACK field.)

\xpart{c} \code{error} stays high until the next frame begins, while \code{tx_done} and
\code{rx_valid} are one-cycle pulses. Explain what would go wrong for the parallel class's C++ driver
if \code{error} were a pulse like the other two.
